%\ifodd\value{page}\hbox{}\newpage\fi
\chapter*{Śrī Lalitā Sahasranāma — Śloka 17}
\addcontentsline{toc}{chapter}{Śloka 17 - Nomi 34,35}

\begin{sloka}
{\sanskritfont
\begin{tabular}{c}
कामेशज्ञातसौभाग्यमार्दवोरुद्वयान्विता ।\\
माणिक्यमुकुटाकारजानुद्वयविराजिता ॥१७॥
\end{tabular}

\vspace{1.5em}

{\middlesize \iastfont
kāmeśa-jñāta-saubhāgya-mārda-voru-dvayānvitā |\\
māṇikya-mukuṭākāra-jānu-dvaya-virājitā ||17||}

\vspace{1em}

{\middlesize
\anvaya{%
कामेशज्ञातसौभाग्यमार्दवोरुद्वयान्विता ।
माणिक्यमुकुटाकारजानुद्वयविराजिता ।
}
}

\middlesize \iastfont \itmean{%
«Dotata di cosce la cui fortuna sponsale e pienezza amorosa,
insieme alla loro tenera morbidezza, sono note a \textit{Kāmeśa};
le cui ginocchia risplendono come corone di rubino.»
}

}
\end{sloka}

\wordmeaning{%
\wordline{कामेश-ज्ञात-सौभाग्य-मार्दव-उरु-द्वय-अन्विता}
{kāmeśa-jñāta-saubhāgya-\\mārda-voru-dvayānvitā}
{ dotata (\textit{anvitā}) di due cosce (\textit{uru-dvaya}) la cui \textit{saubhāgya} indica bellezza e fortuna, ma anche pienezza sponsale e felicità coniugale, e la cui \textit{mārda} esprime morbidezza, delicatezza tattile e tenerezza affettiva, conosciute (\textit{jñāta}) da \textit{Kāmeśa}; }%

\wordline{माणिक्य-मुकुट-आकार-जानु-द्वय-विराजिता}
{māṇikya-mukuṭākāra-jānu-\\dvaya-virājitā}
{ i cui due ginocchi (\textit{jānu-dvaya}) risplendono (\textit{virājitā}) come (\textit{ākāra}) corone (\textit{mukuṭa}) di rubino (\textit{māṇikya}); }%
}

\textit{Śloka} 17 prosegue la contemplazione discendente soffermandosi sulle cosce e sui ginocchi. Le cosce sono descritte come portatrici di una bellezza e di una pienezza sponsale riconosciute da \textit{Kāmeśa}: non una conoscenza esterna o visiva, ma un’intimità consapevole, in cui la fortuna amorosa e la tenerezza della forma sono pienamente condivise.

Nel secondo \textit{pāda}, i ginocchi sono paragonati a corone di rubino. L’immagine non rimanda a un ornamento applicato, ma a una regalità intrinseca della forma stessa. In alcune letture, la forma “a corona” richiama anche le ginocchia che si piegano nella prostrazione: la regalità della Dea si manifesta così non solo nello stare eretta, ma anche nel chinarsi in amore verso il mondo.

Il termine \textit{saubhāgya} non indica una fortuna generica né un semplice attributo estetico. Nella tradizione indiana, esso designa una condizione di pienezza relazionale e sponsale: essere nel legame giusto, amato e riconosciuto, in una forma di compimento che unisce bellezza, desiderabilità e armonia. In questo \textit{śloka}, il \textit{saubhāgya} delle cosce di \textit{Lalitā} non è esibito né visibile, ma è detto “conosciuto” (\textit{jñāta}) da \textit{Kāmeśa}. La qualità sponsale non appartiene allo sguardo esterno, ma all’intimità condivisa: è una pienezza che esiste solo nella relazione, non come ornamento del corpo, ma come verità del legame.

Questo \textit{śloka} contiene due \textit{nāma}.  
\textbf{\textit{Kāmeśa-jñāta-saubhāgya-mārda-voru-dvayānvitā}} presenta le cosce come luogo di pienezza sponsale e tenerezza conosciuta nell’intimità divina.  
\textbf{\textit{Māṇikya-mukuṭākāra-jānu-dvaya-virājitā}} raffigura i ginocchi come segni di regalità luminosa, capaci anche di piegarsi senza perdere sovranità.
